\documentclass{article}

\begin{document}



\section{Lecture 1: Neuromorphic Engineering}
\subsection{organisation}

potentially can get credits for this course but I need to contact the board of examiners asap.

%could you add me to the brightspace page manually? (so I should contact the BoE to get this course registered as part of my degree????)


%has the paper of the full brain mapping affected your research/research projects?



%could you explain the difference between weak and strong inversion, and also the effects of the regieme and also what causes a transistor to operate under these modalities?
%cmos structure??+
%cadence little arrow on transistors?
%why always real time temporal dynamics?? 60-70ms



exam will involve knowing the equations for a first questions


activty = mean firing rate.

%just so that I can play around, are memristors available on cadence?
%if yes could I maybe play around with crossbar network of memristors?? It would be interesting


%for understanding of the brain is it also reasonable to attempt non realtime temporal dynamics?


%Derivations of the transfer functions??? and also derivation for the early voltage (rather than just an empirical motivation???)



\section{Tutorial one}
\subsection{1}

insert drawing of nmos
\subsection{2}


equation describing transistor in subthreshold region. in terms of $I_r$ and $I_f$ 

I = If - Ir
where $U_t$ is the thermal voltage

\(I_{ds}= I_0 e^{\frac{\kappa V_g}{V_c}}\left( e^{\frac{-V_s}{U_t}}-e^{\frac{-V_d}{U_t}}\right)\)
(should have this equation memorised. full derivation can be found in the book (where???)) 


\subsection{3}
what is the effect of changing the source and draing voltage on the forward and reverse currents. 

1. increasing \(V_s\)  $\rightarrow$ decrease $I_ds$ 
2. $V_d$ increacing $I_ds \rightarrow $ increase $I_ds$


\begin{equation}
    \frac{I_f}{I_r} = \frac{e^{-V_s/U_t}}{e^{-V_d/U_t}}= e^{V_{ds}/U_t}
\end{equation}

what happens in the ohmic and saturation regions? $\rightarrow$

the ohmic region is that $V_{ds} (lessthanorequalto) U_t $
saturation region: $V_{ds} > U_t$


\subsection{4}
$V_{ds} V_d - V_s $

thus in the ohmic region
\(\frac{I_f}{I_r} \approx 1\)
and in the saturation region
\(\frac{I_F}{I_r}>>1\)


\subsection{5}
how are the ohmic and saturation regions defined in terms of $V_ds$? write down approximate current equations.

\begin{equation}
    I_{ds}= I_0 e^{\frac{\kappa V_g - V_s}{U_t} (1-e ^{-V_ds / U_t})}
\end{equation}

Sat: \(V_ds > nU_t\) n??? not sure what he is writing on the board there. maybe its kappa

\begin{equation}
    I_ds \approx I_0 e^{\frac{\kappa V_g - V_s}{U_t}}
\end{equation}

ohmic region
\(e^x = 1+x+....\)
\begin{equation}
    I_{ds}= I_0e^{\frac{\kappa V_g - V_s}{U_t} \left(\frac{V_{ds}}{U_t}\right)}
\end{equation}

Vs(Vg)
\begin{equation}
    \frac{I_ds}{I_0}=e^{\frac{\kappa V_g-V_s}{U_t}}
\end{equation}

\begin{equation}
U_t ln (\frac{I}{I_0}) = \kappa V_g - V_s
\end{equation}

\begin{equation}
    V_s = \kappa V_g - U_t \ln \left(\frac{I_{ds}}{I_0}\right)
\end{equation}


\subsection{...}
insert image of pmos.


for both nmos and pmos little arrow is on the source side.

for pmos ground is the drain 

for nmos ground is the source



\end{document}